\section{Data}
\label{sec:data}

\subsection{Data Source}

We use the Encuesta Nacional de Hogares (ENAHO) Module 500 (Employment), Peru's primary household labor force survey conducted by INEI. Our panel covers 2020--2024, with 348,505 individuals and 7,337 variables in wide format. The data is reshaped to long format, yielding 1,742,525 individual-year observations.

\subsection{Sample}

We restrict to working-age (15--65) employed individuals (\texttt{ocu500} = 1), yielding 262,733 observations across 214,413 unique individuals. ENAHO uses a rotating panel design where households are tracked for approximately two consecutive years before replacement; only 0.5\% of individuals are observed in all five waves.

\subsection{Outcome: Informality}

Our primary informality measure is the absence of social security coverage (\texttt{p511a} = 7, ``no insurance through employer''). This captures 19.5\% of employed workers in 2020, rising sharply to 34.5\% in 2021 and stabilizing around 34\% through 2024. We construct two alternative measures for robustness: (i) absence of a written contract (\texttt{p517} $\in \{5, 6\}$) and (ii) employment in a small firm with fewer than 5 workers (\texttt{p510} $\in \{1, 2\}$).

\subsection{Treatment: Teleworkability}

We measure treatment intensity using the \citet{dingel2020} teleworkability index, mapped to ENAHO via a crosswalk from BLS SOC codes to ISCO-08 2-digit occupation codes. The crosswalk aggregates D\&N's 764 occupation-level teleworkability scores to 43 ISCO-08 sub-major groups, then merges with ENAHO's \texttt{p505r4} occupation variable (which follows the ISCO-08 classification). The match rate is 100\% for employed workers.

Our primary treatment measure is the \textit{continuous} teleworkability score (0--1). As a secondary specification, we define a binary indicator: \textbf{contact-intensive} = teleworkability score $< 0.20$. This captures 74.6\% of the employed sample, reflecting Peru's concentration of workers in occupations requiring physical presence (agriculture, construction, personal services, manufacturing). Only 8.2\% of workers are in highly teleworkable occupations (score $\geq 0.50$).

Table~\ref{tab:descriptive} compares contact-intensive and teleworkable workers. Contact-intensive workers have higher baseline informality rates and lower incomes, consistent with the well-documented correlation between occupation type and labor formality in Peru.

\begin{table}[htbp]
\centering
\caption{Descriptive Statistics by Teleworkability}
\label{tab:descriptive}
\begin{tabular}{lccccc}
\toprule
 & \multicolumn{2}{c}{Contact-Intensive} & \multicolumn{2}{c}{Teleworkable} & \\
\cmidrule(lr){2-3} \cmidrule(lr){4-5}
Variable & Mean & N & Mean & N & Diff \\
\midrule
Informality rate & 0.358 & 196,069 & 0.192 & 66,664 & 0.166 \\
Age & 40 & 196,069 & 40 & 66,664 & 0.570 \\
Female & 0.416 & 196,069 & 0.610 & 66,664 & -0.194 \\
Rural & 0.442 & 196,069 & 0.123 & 66,664 & 0.319 \\
Monthly income (soles) & 15440 & 78,180 & 25408 & 38,777 & -9968 \\
\bottomrule
\end{tabular}
\begin{tablenotes}\small
\item Notes: Contact-intensive = teleworkability score $<$ 0.20. Teleworkable = score $\geq$ 0.20. Working-age employed sample (15--65), 2020--2024.
\end{tablenotes}
\end{table}


\subsection{Survey Design}

All estimates use ENAHO's survey expansion factor (\texttt{facpob07}) as sampling weights. Standard errors are clustered at the ISCO-08 2-digit occupation level (40 clusters), the level at which treatment is assigned. We discuss the implications of the rotating panel design for within-individual estimation in Section~\ref{sec:robustness}.
